% Os tempos de acesso da hierarquia, em escala logarítmica, com a tradução
% para uma escala humana.
%
% A escala é logarítmica porque a diferença entre os extremos é de nove ordens
% de grandeza: em escala linear, os cinco primeiros níveis cairiam todos sobre
% o mesmo ponto. A coluna da direita reescala tudo por um fator único, tomando
% um ciclo de relógio como um segundo, e é o que torna a diferença
% perceptível.
\documentclass[border=5pt]{standalone}
\usepackage{figuras}
\begin{document}
\begin{tikzpicture}

  % --- Eixo logarítmico ---------------------------------------------------
  \draw[ligacao, -{Stealth[length=2.6mm]}] (-0.1, 0) -- (14.6, 0);
  \foreach \x/\rot in {0.87/1 ns, 5.22/1 $\mu$s, 9.57/1 ms, 13.92/1 s} {
    \draw[ligacao] (\x, -0.12) -- (\x, 0.12);
    \node[rotulo peq, below=1mm] at (\x, -0.1) {\rot};
  }
  \node[rotulo peq, anchor=north east] at (14.6, -0.55)
    {tempo de acesso, em escala logarítmica};

  % --- Um nível por linha -------------------------------------------------
  % {x}{y}{cor}{nome}{tempo real}{tempo na escala humana}
  \newcommand{\nivel}[6]{%
    \draw[dash pattern=on 2pt off 2pt, draw=corapagado] (#1, 0.1) -- (#1, #2);
    \node[circle, fill=#3, draw=#3, inner sep=2.4pt] at (#1, #2) {};
    \node[rotulo peq, anchor=east] at (-0.4, #2) {#4};
    \node[rotulo peq, anchor=south west, text=#3] at (#1 + 0.12, #2 + 0.06)
      {#6};
    \node[rotulo peq, anchor=north west] at (#1 + 0.12, #2 - 0.02)
      {\scriptsize #5};
  }

  \nivel{12.73}{6.3}{esborda}{Ida e volta a outro continente}{cerca de 150 ms}{14 anos}
  \nivel{11.02}{5.4}{dadoborda}{Disco magnético}{cerca de 10 ms}{1 ano}
  \nivel{7.68}{4.5}{dadoborda}{Disco de estado sólido}{dezenas de $\mu$s}{2 dias}
  \nivel{3.77}{3.6}{memborda}{Memória principal}{cerca de 100 ns}{5 minutos}
  \nivel{2.49}{2.7}{procborda}{Memória cache, nível 3}{cerca de 13 ns}{40 segundos}
  \nivel{1.04}{1.8}{procborda}{Memória cache, nível 1}{cerca de 1 ns}{4 segundos}
  \nivel{0.17}{0.9}{procborda}{Registrador}{um ciclo de relógio}{1 segundo}

  \node[rotulo peq, anchor=east, font=\footnotesize\bfseries] at (-0.4, 7.1)
    {Onde o valor está};
  \node[rotulo peq, anchor=west, font=\footnotesize\bfseries] at (-0.2, 7.1)
    {Se um ciclo de relógio durasse um segundo};

\end{tikzpicture}
\end{document}
